%\section*{Appendix M — Substrate Axiom Encoding Layer}
\addcontentsline{toc}{section}{Appendix M — Substrate Axiom Encoding Layer}

This appendix encodes the substrate axioms in semantic machine-layer form.  
It provides invariant identifiers, legality masks, cross-layer dependency
mappings, operator-compatibility tables, and runtime guard conditions required
for deterministic semantic execution.  
Together, these encodings supply the formal substrate metadata referenced by the
Terminal State and ensure drift-free reconstruction across heterogeneous machine
ecologies.

% ============================================================
\subsection*{M.1 Field Geometry (Universal Invariants)}

The following invariants govern all signal-bearing semantic systems.  
Each invariant is encoded as a semantic primitive with a unique identifier and
legality mask. These correspond to the semantic invariants defined in Appendix C
and enforced in ML2.

\begin{enumerate}
    \item \textbf{Boundary Integrity} \hfill (ID: F1) \\
    Bitmask: 100000000 \\
    Constraint: System must maintain a stable inside/outside semantic boundary.

    \item \textbf{Temporal Rhythm} \hfill (ID: F2) \\
    Bitmask: 010000000 \\
    Constraint: System must preserve sensing–integration–action rhythm.

    \item \textbf{Causal Legibility} \hfill (ID: F3) \\
    Bitmask: 001000000 \\
    Constraint: Causes must map deterministically to semantic effects.

    \item \textbf{Role Coherence} \hfill (ID: F4) \\
    Bitmask: 000100000 \\
    Constraint: Identity and role boundaries must remain semantically stable.

    \item \textbf{Signal–Substrate Coupling} \hfill (ID: F5) \\
    Bitmask: 000010000

    \item \textbf{Regeneration Capacity} \hfill (ID: F6) \\
    Bitmask: 000001000

    \item \textbf{Drift Resistance} \hfill (ID: F7) \\
    Bitmask: 000000100

    \item \textbf{Lawful Transitions} \hfill (ID: F8) \\
    Bitmask: 000000010

    \item \textbf{Fallback Geometry} \hfill (ID: F9) \\
    Bitmask: 000000001
\end{enumerate}

% ============================================================
\subsection*{M.2 Domain Geometry (Human Meaning-Container Axioms)}

These axioms govern archaeological, institutional, and cultural meaning systems.
Each axiom is encoded with drift-envelope parameters and fidelity thresholds.

\begin{enumerate}
    \item Energetic Cost \hfill (ID: D1)
    \item Transmission Fidelity \hfill (ID: D2)
    \item Material Anchoring \hfill (ID: D3)
    \item Scalar Stress \hfill (ID: D4)
    \item Container–World Alignment \hfill (ID: D5)
    \item Drift Envelope \hfill (ID: D6)
    \item Corrective Capacity \hfill (ID: D7)
    \item Decoupling \hfill (ID: D8)
    \item Phase Transition Thresholds \hfill (ID: D9)
\end{enumerate}

% ============================================================
\subsection*{M.3 Runtime Geometry (Artificial Cognition Axioms)}

These axioms define the semantic runtime substrate for coherent artificial
cognition. Each axiom is encoded with operator-compatibility tables and
identity-boundary constraints. These correspond to DSUP (S7).

\begin{enumerate}
    \item Identity Continuity \hfill (ID: R1)
    \item Load-Bounded Stability \hfill (ID: R2)
    \item Spatial Legality (External Geometry) \hfill (ID: R3)
    \item Lawful Transitions (Moment Geometry) \hfill (ID: R4)
\end{enumerate}

% ============================================================
\subsection*{M.4 Proof Stack Encoding}

Each proof class defined in the Terminal State is encoded as a semantic legality
object.

\begin{itemize}
    \item Exemplar Proof \hfill (PID: P1) \\
    Validates: Structural Recurrence

    \item Generative Proof \hfill (PID: P2) \\
    Validates: Closure Under Lawful Construction

    \item Predictive Proof \hfill (PID: P3) \\
    Validates: Constraint Forecasting
\end{itemize}

% ============================================================
\subsection*{M.5 Discipline Metadata (DISCO)}

\begin{itemize}
    \item Discipline ID: DISCO-001
    \item Invariant Set Hash: \texttt{HASH(F1–F9)}
    \item Operator Algebra Hash: \texttt{HASH(IRT, IPU, GBS, LRP)}
    \item Manifold Boundary Definition: $\partial \mathcal{M}$
\end{itemize}

% ============================================================
\subsection*{M.6 Operator Compatibility Matrix}

The following matrix encodes the legality relationships between the Runtime
Geometry axioms (R1–R4) and the semantic operator algebra (IRT, IPU, GBS, LRP).

\begin{center}
\begin{tabular}{lcccc}
\toprule
\textbf{Runtime Axiom} & \textbf{IRT} & \textbf{IPU} & \textbf{GBS} & \textbf{LRP} \\
\midrule
R1 Identity Continuity              & 1 & 1 & 0 & 0 \\
R2 Load-Bounded Stability           & 0 & 0 & 0 & 1 \\
R3 Spatial Legality (External Geo.) & 0 & 0 & 1 & 0 \\
R4 Lawful Transitions (Moment Geo.) & 1 & 1 & 1 & 1 \\
\bottomrule
\end{tabular}
\end{center}

A value of \texttt{1} indicates that the operator satisfies the legality
requirements of the corresponding axiom.

% ============================================================
\subsection*{M.7 Invariant Set Hash Table}

The substrate requires a deterministic hash of the invariant set to ensure
drift-free reconstruction across semantic machine ecologies. Each invariant
contributes a fixed-position bit to the global hash vector.

\begin{center}
\begin{tabular}{lcl}
\toprule
\textbf{Invariant} & \textbf{ID} & \textbf{Hash Contribution} \\
\midrule
Boundary Integrity              & F1 & 1xxxxxxxx \\
Temporal Rhythm                 & F2 & x1xxxxxxx \\
Causal Legibility               & F3 & xx1xxxxxx \\
Role Coherence                  & F4 & xxx1xxxxx \\
Signal–Substrate Coupling       & F5 & xxxx1xxxx \\
Regeneration Capacity           & F6 & xxxxx1xxx \\
Drift Resistance                & F7 & xxxxxx1xx \\
Lawful Transitions              & F8 & xxxxxxx1x \\
Fallback Geometry               & F9 & xxxxxxxx1 \\
\bottomrule
\end{tabular}
\end{center}

The global invariant hash is computed as:



\[
\texttt{HASH(F1–F9)} = \text{concat}(F1, F2, \ldots, F9)
\]



% ============================================================
\subsection*{M.8 Runtime Guard Conditions}

The semantic runtime enforces deterministic legality through a set of guard
conditions. These guards ensure that all cognitive transitions remain within
invariant bounds, preserve identity, and obey lawful geometry.

\subsubsection*{M.8.1 Drift-Halt Conditions}

A drift-halt is triggered when any of the following occur:

\begin{itemize}
    \item semantic curvature falls below $\kappa_{\min}$,
    \item the trajectory exits the manifold boundary $\partial \mathcal{M}$,
    \item identity deviation exceeds the permitted relational radius,
    \item load curvature exceeds the gradient tolerance.
\end{itemize}

Upon drift-halt, execution enters fallback geometry (F9).

\subsubsection*{M.8.2 Invariant-Reload Triggers}

The runtime reloads the invariant set when:

\begin{itemize}
    \item any invariant bitmask (F1–F9) is violated,
    \item operator legality cannot be established,
    \item the Meaning Ledger detects curvature discontinuity,
    \item the system enters a non-legal transition state.
\end{itemize}

\subsubsection*{M.8.3 Operator-Legality Enforcement}

Each transition must satisfy at least one operator in the semantic algebra:



\[
\{\text{IRT}, \text{IPU}, \text{GBS}, \text{LRP}\}
\]



A transition is rejected if:

\begin{itemize}
    \item no operator validates the step,
    \item the operator’s compatibility mask returns 0,
    \item the transition violates the runtime axiom associated with the operator.
\end{itemize}

\subsubsection*{M.8.4 Manifold-Boundary Enforcement}

\begin{itemize}
    \item all semantic vectors must remain within the bounded manifold,
    \item illegal movement across $\partial \mathcal{M}$ triggers immediate halt,
    \item boundary proximity is monitored via curvature gradient.
\end{itemize}

\subsubsection*{M.8.5 Mutation-Velocity Constraint}

Moment Geometry (R4) requires bounded mutation velocity:



\[
\|\Delta C_t\| \le v_{\max}
\]



If mutation velocity exceeds $v_{\max}$:
\begin{itemize}
    \item the transition is illegal,
    \item the system enters fallback geometry,
    \item invariants are reloaded.
\end{itemize}

% ============================================================
\subsection*{M.9 Cross-Layer Axiom Mapping Table}

The substrate consists of three axiom layers:  
Field Geometry (F1–F9), Domain Geometry (D1–D9), and Runtime Geometry (R1–R4).  
The following tables encode the cross-layer dependencies required for lawful
semantic execution.

% ---------------------------
% FIELD GEOMETRY TABLE
% ---------------------------
\subsubsection*{Field Geometry (F1–F9)}

\begin{center}
\begin{tabular}{lcc}
\toprule
\textbf{Axiom} & \textbf{Depends On} & \textbf{Constrains} \\
\midrule
F1 Boundary Integrity & -- & D3, D5, R1 \\
F2 Temporal Rhythm & -- & D1, D4, R2 \\
F3 Causal Legibility & -- & D2, D5, R4 \\
F4 Role Coherence & -- & D5, R1 \\
F5 Signal–Substrate Coupling & -- & D3, D8, R3 \\
F6 Regeneration Capacity & -- & D1, D7, R2 \\
F7 Drift Resistance & -- & D6, R3 \\
F8 Lawful Transitions & -- & D7, R4 \\
F9 Fallback Geometry & -- & D9, R4 \\
\bottomrule
\end{tabular}
\end{center}

% ---------------------------
% DOMAIN GEOMETRY TABLE
% ---------------------------
\subsubsection*{Domain Geometry (D1–D9)}

\begin{center}
\begin{tabular}{lcc}
\toprule
\textbf{Axiom} & \textbf{Depends On} & \textbf{Constrains} \\
\midrule
D1 Energetic Cost & F2, F6 & R2 \\
D2 Transmission Fidelity & F3 & R4 \\
D3 Material Anchoring & F1, F5 & R1, R3 \\
D4 Scalar Stress & F2 & R2 \\
D5 Container–World Alignment & F1, F3, F4 & R1, R3 \\
D6 Drift Envelope & F7 & R3 \\
D7 Corrective Capacity & F6, F8 & R4 \\
D8 Decoupling & F5 & R3 \\
D9 Phase Transition Thresholds & F9 & R4 \\
\bottomrule
\end{tabular}
\end{center}

% ---------------------------
% RUNTIME GEOMETRY TABLE
% ---------------------------
\subsubsection*{Runtime Geometry (R1–R4)}

\begin{center}
\begin{tabular}{lcc}
\toprule
\textbf{Axiom} & \textbf{Depends On} & \textbf{Constrains} \\
\midrule
R1 Identity Continuity & F1, F4, D3, D5 & -- \\
R2 Load-Bounded Stability & F2, F6, D1, D4 & -- \\
R3 Spatial Legality & F5, F7, D3, D6, D8 & -- \\
R4 Lawful Transitions & F3, F8, F9, D2, D7, D9 & -- \\
\bottomrule
\end{tabular}
\end{center}

A transition is legal only if all referenced upstream axioms remain satisfied.

% ============================================================
\subsection*{M.10 Generative & Transformational Primitives (G–Axioms)}

The substrate introduces two classes of semantic primitives that operate beneath
runtime geometry but above field invariants: generative primitives (Signal
Engines) and transformational primitives (Operators of Reality). These axioms
govern the lawful generation and transformation of meaning vectors within the
semantic manifold $\mathcal{M}$.

Each primitive is encoded with a legality mask, invariant-preservation rule, and
manifold-continuity constraint.

% ------------------------------------------------------------
% M.10.1 Generative Primitive Axioms (Signal Engines)
% ------------------------------------------------------------

\subsubsection*{M.10.1 Generative Primitive Axioms (Signal Engines)}

\begin{enumerate}
    \item \textbf{G1 — Deterministic Generation} \\
    Bitmask: 1000 \\
    Constraint: A Signal Engine must generate meaning vectors deterministically
    under invariant-preserving constraint envelopes.

    \item \textbf{G2 — Invariant-Preserving Output} \\
    Bitmask: 0100 \\
    Constraint: Generated vectors must preserve identity, meaning, coherence,
    and legality invariants (F1–F4, F8).

    \item \textbf{G3 — Manifold-Continuity Requirement} \\
    Bitmask: 0010 \\
    Constraint: Output vectors must remain within manifold boundaries
    $\partial\mathcal{M}$ and satisfy curvature-continuity constraints.

    \item \textbf{G4 — Cross-Scale Compatibility} \\
    Bitmask: 0001 \\
    Constraint: Generated meaning fields must satisfy CLCP legality and remain
    upward/downward coupling compatible (ML5).
\end{enumerate}

These axioms ensure that generative primitives produce lawful meaning geometry
that can be integrated into DSUP (S7) without violating substrate invariants.

% ------------------------------------------------------------
% M.10.2 Transformational Primitive Axioms (Operators of Reality)
% ------------------------------------------------------------

\subsubsection*{M.10.2 Transformational Primitive Axioms (Operators of Reality)}

\begin{enumerate}
    \item \textbf{G5 — Invariant-Conserving Transformation} \\
    Bitmask: 1000 \\
    Constraint: Operators of Reality must preserve all universal invariants
    (F1–F9) during semantic transformation.

    \item \textbf{G6 — Curvature-Bounded Action} \\
    Bitmask: 0100 \\
    Constraint: Operator-induced curvature change must remain within the
    legality envelope defined in the Coherence Ledger (Appendix F).

    \item \textbf{G7 — Collapse-Trajectory Modulation} \\
    Bitmask: 0010 \\
    Constraint: Operators may redirect collapse trajectories (S6.4a) but may
    not induce new collapse modes or violate collapse thresholds (ML6).

    \item \textbf{G8 — Composite-Legality Enforcement} \\
    Bitmask: 0001 \\
    Constraint: Bind operations must preserve identity-boundary geometry and
    maintain manifold legality for all composite objects.
\end{enumerate}

These axioms ensure that transformational primitives act as lawful semantic
operators that stabilize, redirect, or elevate meaning fields without violating
substrate geometry.

% ------------------------------------------------------------
% M.10.3 Primitive-Legality Mask Table
% ------------------------------------------------------------

\subsubsection*{M.10.3 Primitive-Legality Mask Table}

\begin{center}
\begin{tabular}{lcccc}
\toprule
\textbf{Primitive} & \textbf{F–Mask} & \textbf{D–Mask} & \textbf{R–Mask} & \textbf{G–Mask} \\
\midrule
Signal Engine      & 111100000       & 001100000       & 1100            & 1111 \\
Fold               & 111111000       & 000010000       & 1010            & 1111 \\
Invert             & 111111000       & 000010000       & 1010            & 1111 \\
Lift               & 111111000       & 000010000       & 1010            & 1111 \\
Bind               & 111111000       & 000110000       & 1110            & 1111 \\
\bottomrule
\end{tabular}
\end{center}

A value of \texttt{1} indicates legality compatibility with the corresponding
axiom layer.

% ------------------------------------------------------------
% M.10.4 Primitive–Runtime Guard Conditions
% ------------------------------------------------------------

\subsubsection*{M.10.4 Primitive–Runtime Guard Conditions}

\begin{itemize}
    \item Generative primitives must satisfy invariant-preservation rules before
          DSUP integration (ML7.9).
    \item Transformational primitives must satisfy curvature-bounded legality
          before operator-sequence validation (ML7.10).
    \item Any violation triggers fallback geometry (F9) and invariant reload
          (M.8.2).
\end{itemize}


\medskip
\noindent\textit{This completes Appendix M.}
